% Self-contained derivations for the diagnostic statements in the main text.
\section{Proofs and supporting counterexamples}
\label{app:proofs}

\subsection{Cue-channel ordering and readout effects}

\paragraph{Proof of \cref{prop:affine}.}
Unrolling the switched-affine recurrence gives
\begin{equation}
 s_H=\Phi_hE(u)+k_h,\qquad
 \Phi_h=A_{x_{H-1}}\cdots A_{x_0},\qquad
 k_h=\sum_{j=0}^{H-1}
       (A_{x_{H-1}}\cdots A_{x_{j+1}})c_{x_j},
 \label{eq:unrolled-switched}
\end{equation}
where an empty product is the identity. Subtraction cancels $k_h$,
differentiation gives $D_us_H=\Phi_hDE(u)$, and covariance is invariant to
translation. Fixed $A$ yields $\Phi_h=A^H$. Pairwise commutativity yields
$\Phi_h=\prod_\sigma A_\sigma^{n_\sigma}$, so equal counts give equal products.
All covariance claims require a common cue law with finite second moments.
For fixed $A$ and cue-independent additive noise,
$s_H=A^HE(u)+k_h+\sum_jA^{H-1-j}\eta_j$. A common joint noise law gives the
same propagated law; if the noise is independent of the cue, its covariance
adds to the displayed signal covariance. With independently drawn episode
noise, individual cue differences also contain noise differences. With
switched matrices, even independent identically distributed process noise can
produce unequal terminal covariance: two scalar transitions $a,b$, with
independent variance-$\tau^2$ noise after each update, give terminal-noise
variances $(1+b^2)\tau^2$ in order $ab$ and $(1+a^2)\tau^2$ in order $ba$,
although both cue products equal $ab$. \qed

\paragraph{Proof of \cref{cor:offset}.}
Subtracting $k_h$ makes the joint law of state and target identical between
bodies. Composing any decoder with this translation gives the same risk under
any fixed loss; the affine-decoder class is also closed under translations
when it includes an intercept. This is an equality of attainable optimal
risks, not of finite-sample fitted probes. Expanding
$\norm{e+\Delta}^2-\norm{e}^2$ proves \cref{eq:offset-gap}.
If bodies vary across paired samples, the general expression is
$2\E[e^\top\Delta]+\E\norm{\Delta}^2$; the factorization in the main text
requires a deterministic shift. The required compensation is body-specific:
a single intercept per arm need not suffice for a mixture of bodies. \qed

\begin{remark}[Noncommutativity is insufficient]
\label{prop:noncommuting}
If equal-count products differ, at least one matrix pair fails to commute,
by the contrapositive of \cref{prop:affine}. The converse is false.
Take $A=\diag(1,-1)$ and
$B=\left(\begin{smallmatrix}0&1\\1&0\end{smallmatrix}\right)$.
Then $AB=-BA$, but $(AB)^2=(BA)^2=-I$.
Moreover, the unequal products $AB$ and $BA$ are orthogonal, so their cue
channels have identical optimal recovery risk with independent isotropic Gaussian terminal noise
and unrestricted decoding: rotating the observation converts either into the
same channel $u+\varepsilon$. Thus a changed product is itself insufficient
for a changed recovery risk.
\end{remark}

\begin{lemma}[Adjacent-swap bound]
\label{lem:swaps}
Let two words of length $H\ge2$ have the same symbol counts and be connected
by $K$ adjacent swaps. If $\norm{A_\sigma}_2\le L$ with $L>0$ and
$\kappa=\max_{\sigma,\tau}\norm{A_\sigma A_\tau-A_\tau A_\sigma}_2$, then
\[
 \norm{\Phi_h-\Phi_{h'}}_2\le K L^{H-2}\kappa.
\]
\end{lemma}
\begin{proof}
For a single swap, factor the difference into a common left product, the
commutator, and a common right product. The common products contain $H-2$
factors in total, so submultiplicativity bounds the difference by
$L^{H-2}\kappa$. Sum the $K$ differences along the swap sequence and use the
triangle inequality. Equal-count words admit such a sequence by successive
adjacent transpositions. The corresponding change in a cue difference is at
most $KL^{H-2}\kappa\norm{E(u)-E(v)}$. This standard telescoping bound omits
affine offsets and process-noise propagation, and has no lower bound or sign.
\end{proof}

\subsection{Covariance geometry and noisy recovery}

\begin{proposition}[Sampling ceilings and coordinate dependence]
\label{prop:ceilings}\label{prop:pr}
A centered cloud of $N$ states in $\R^n$ has covariance rank at most
$\min(N-1,n)$. A time-only sample of an exactly $P$-periodic orbit has rank at
most $\min(P-1,n)$. For nonzero covariance $C\succeq0$,
$\PR(C)=(\tr C)^2/\tr(C^2)$ is invariant to positive scalar multiplication
and orthogonal coordinate changes, but not to anisotropic invertible scaling.
Writing $p_i=\lambda_i/\tr C$ on the positive eigenvalues,
\[
 1\le\PR(C)=e^{H_2(p)}\le e^{H_1(p)}\le\rank C.
\]
The upper three quantities coincide exactly for equal positive eigenvalues;
equality also with one requires rank one. For zero covariance, normalized
dimension is undefined.
\end{proposition}
\begin{proof}
Centered sample rows sum to zero. An orbit sample additionally lies in the
affine span of at most $P$ points. These give the rank bounds. Scaling and
orthogonal invariance follow directly from traces. In contrast,
$C=I_2$ transformed by $\diag(1,\epsilon)$ has
$\PR=(1+\epsilon^2)^2/(1+\epsilon^4)\to1$ as $\epsilon\downarrow0$,
although the transformation is invertible for every $\epsilon>0$.
Finally, $\PR=1/\sum_i p_i^2$. Jensen's inequality gives
$-\sum_i p_i\log p_i\ge-\log\sum_i p_i^2$, and the entropy maximum on a
support of size $r$ is $\log r$, attained exactly at the uniform law.
\end{proof}

\begin{lemma}[Linear recovery channel]
\label{prop:channel}
For independent zero-mean $u,\varepsilon$ with covariances
$\Lambda\succ0,\Sigma\succ0$, the best linear estimator of $u$ from
$y=Gu+\varepsilon$ has squared error
$\tr[(\Lambda^{-1}+G^\top\Sigma^{-1}G)^{-1}]$.
For Gaussian $u,\varepsilon$, this is also the unrestricted Bayes risk.
\end{lemma}
\begin{proof}
The best linear estimator is
$\Lambda G^\top(G\Lambda G^\top+\Sigma)^{-1}y$.
Its error covariance is
$\Lambda-\Lambda G^\top(G\Lambda G^\top+\Sigma)^{-1}G\Lambda$;
the Woodbury identity yields the displayed expression. Gaussian conditioning
makes the same estimator the conditional mean. Setting
$G=a^HI_d$, $\Lambda=I_d$, and $\Sigma=\sigma^2I_d$ gives
\cref{eq:attenuation-risk}. For non-Gaussian cues, a nonlinear decoder can
improve on the linear risk.
\end{proof}

For the complementary nuisance-compression example in the main text, put
$q=a^{2H}$. The covariance of $(u,a^Hv)$ is
$\diag(I_d,qI_k)$ and its participation ratio is
$(d+kq)^2/(d+kq^2)$. Its derivative with respect to $q\in(0,1)$ is
$2dk(d+kq)(1-q)/(d+kq^2)^2>0$ for $d,k>0$, so participation ratio decreases
with $H$ while the readout selecting $u$ remains exact.

\subsection{Categorical basins and continuous memory}

\paragraph{Proof of \cref{prop:basins}.}
The shifted coordinate satisfies $r_{t+1}=g(r_t)$, with
$g(\pm1)=\pm1$ and $g'(\pm1)=0$. At $r=1+\delta$,
$g(r)-1=-1.5\delta^2-0.5\delta^3$; odd symmetry gives the same magnitude
bound at $-1$. For $|\delta|\le0.05$, the next error is at most
$0.0038125<0.05$, proving invariance, and
$|g'(r)|\le1.5(2\cdot0.05+0.05^2)=0.15375<1$, proving contraction.
At each phase the two tubes are separated by at least $1.9$;
$|\phi_t|=0.2$ ensures that their signs also agree with their labels.
The period-two tangent product along either reference response is zero.
This proves retention for all times within the stated tubes, not global
attraction to those tubes. \qed

\paragraph{Proof of \cref{prop:continuous}.}
Differentiating $\mathcal D(E(u))=u$ yields
$D\mathcal D(E(u))DE(u)=I_d$, so $DE(u)$ has rank $d$.
Differentiating $\mathcal P(E(u))=E(u)$ yields
$D\mathcal P(E(u))DE(u)=DE(u)$. Each of the $d$ independent columns of
$DE(u)$ is therefore a unit-eigenvalue eigenvector. Both identities must hold
on an open set. A small average reconstruction error cannot replace them:
in one dimension, $E_\epsilon(u)=u+\epsilon\sin(u/\epsilon)$ and
$\mathcal D(s)=s$ have pointwise squared error at most $\epsilon^2$,
yet $E_\epsilon'$ vanishes wherever $\cos(u/\epsilon)=-1$.
Approximate competence alone therefore gives no positive local gain bound. \qed

\subsection{Periodic response, transient gain, and reinjection}

\begin{remark}[A forced oscillation needs no neutral multiplier]
\label{prop:scalar}
For $s_{t+1}=as_t+bx_t$ with $0<|a|<1$, $b\ne0$, and $T$-periodic $x$,
\[
 s_t^*=\sum_{k=0}^{\infty}a^k b x_{t-1-k}
\]
is the unique periodic solution and unique bounded bi-infinite solution.
Absolute convergence and an index shift verify the recurrence; input
periodicity verifies $s_{t+T}^*=s_t^*$. Every forward solution has
$s_t-s_t^*=a^t(s_0-s_0^*)$, so the response attracts globally and its period
multiplier is $a^T$. For an input harmonic $e^{i\omega t}$, the response gain
is $b/(e^{i\omega}-a)$, which is nonzero. An input harmonic can therefore
persist while initial-state sensitivity vanishes. Minimal period $T$ does
not itself guarantee a nonzero first harmonic at frequency $1/T$.
\end{remark}

\begin{remark}[Stable eigenvalues do not exclude transient gain]
\label{prop:nonnormal}
For $M=\left(\begin{smallmatrix}a&g\\0&a\end{smallmatrix}\right)$ with
$|a|<1<g$, the spectral radius is $|a|<1$, yet
$\norm{M}_2\ge\norm{Me_2}_2=\sqrt{g^2+a^2}>1$.
Stable period multipliers concern asymptotic powers of the period map;
finite-window and within-period gains require their own singular-value
measurements. They are not implied to decrease monotonically.
\end{remark}

\begin{remark}[Reinjected input is a different memory problem]
\label{prop:reinjection}
For a $C^1$ depth loop $z_{\ell+1}=F(z_\ell,x)$, differentiation gives
$S_{\ell+1}=J_\ell S_\ell+B_\ell$ with $J_\ell=\partial_zF$ and
$B_\ell=\partial_xF$. At a fixed point with $\rho(J^*)<1$, the implicit
function theorem gives local fixed-point sensitivity
$S^*=(I-J^*)^{-1}B^*$ \citep{deep_equilibrium}.
For $z_{\ell+1}=az_\ell+(1-a)x$ with $|a|<1$, the explicit solution is
$z_\ell=a^\ell z_0+(1-a^\ell)x$: initial-state influence vanishes while
input sensitivity approaches one. This does not establish retention of a
cue absent from subsequent inputs.
\end{remark}
